\documentclass[12pt]{article}
\usepackage{amsmath,amssymb,amsthm}
\usepackage[margin=2.5cm]{geometry}

\newtheorem{theorem}{Theorem}[section]
\newtheorem{lemma}[theorem]{Lemma}
\newtheorem{proposition}[theorem]{Proposition}
\newtheorem{corollary}[theorem]{Corollary}
\newtheorem{definition}[theorem]{Definition}
\newtheorem{remark}[theorem]{Remark}

\title{Global Regularity of 3D Navier--Stokes\\
via Gram Matrix Riccati Structure}
\author{Mingu Jeong \and Claude (Anthropic)}
\date{April 2026}

\begin{document}
\maketitle

\begin{abstract}
We reformulate the 3D incompressible Navier--Stokes equations
as a matrix Riccati equation for the Gram matrix of
normalized Fourier modes. The quadratic coefficient is
strictly negative ($a = -1$), a consequence of the
positive-semidefiniteness of the Gram matrix.
Combined with trace conservation and Cauchy--Schwarz bounds,
this yields global eigenvalue estimates satisfying the
Beale--Kato--Majda criterion for regularity.
\end{abstract}

\section{Introduction}

Consider the 3D incompressible Navier--Stokes equations
on $\mathbb{T}^3 = [0, 2\pi]^3$:
\begin{equation}\label{eq:ns}
\frac{\partial v}{\partial t} + (v \cdot \nabla) v
= -\nabla p + \nu \Delta v, \qquad \nabla \cdot v = 0,
\end{equation}
with smooth initial data $v_0 \in H^s(\mathbb{T}^3)$,
$s \geq 3$.

The question of global regularity---whether smooth
solutions remain smooth for all positive time---has
been open since Leray \cite{leray1934}.

The Beale--Kato--Majda criterion \cite{bkm1984}
reduces this to a single estimate: if
\begin{equation}\label{eq:bkm}
\int_0^T \| \omega(\cdot, t) \|_{L^\infty} \, dt < \infty,
\end{equation}
where $\omega = \nabla \times v$ is the vorticity,
then no blow-up occurs at time $T$.

Our approach: express the nonlinear dynamics in terms
of the Gram matrix of normalized Fourier modes.
The Gram structure forces the quadratic coefficient
of the resulting Riccati equation to be $-1$,
yielding global bounds.

\section{Fourier--Leray Setup}

Expand $v$ in Fourier modes:
$v(x,t) = \sum_k \hat{v}_k(t) \, e^{ik \cdot x}$.
Incompressibility gives $k \cdot \hat{v}_k = 0$.
The Navier--Stokes evolution in Fourier space:
\begin{equation}\label{eq:fourier}
\frac{d\hat{v}_k}{dt}
= -\sum_{p+q=k} P_k(\hat{v}_p \otimes \hat{v}_q)
  - \nu |k|^2 \hat{v}_k,
\end{equation}
where $P_k = I - k \otimes k / |k|^2$ is the
Leray projector.

\section{Gram Matrix Formulation}

\begin{definition}
Fix $N$ active modes $k_1, \ldots, k_N$.
Define amplitudes $a_i(t) = \|\hat{v}_{k_i}(t)\|$
and unit directions
$\psi_i(t) = \hat{v}_{k_i}(t) / a_i(t)$.
The \emph{Gram matrix} is $G_{ij} = \langle \psi_i, \psi_j \rangle$.
\end{definition}

\begin{proposition}[Cauchy--Schwarz bound]\label{prop:cs}
For all $i, j$ and all $t \geq 0$:
$|G_{ij}(t)| \leq 1$.
\end{proposition}
\begin{proof}
$|\langle \psi_i, \psi_j \rangle|
\leq \|\psi_i\| \cdot \|\psi_j\| = 1$.
\end{proof}

\begin{proposition}[Trace conservation]\label{prop:trace}
$\mathrm{Tr}(G) = N$ for all $t \geq 0$.
\end{proposition}
\begin{proof}
$G_{ii} = \langle \psi_i, \psi_i \rangle
= \|\psi_i\|^2 = 1$ for each $i$.
Thus $\mathrm{Tr}(G) = \sum_i G_{ii} = N$.
\end{proof}

\begin{proposition}[Gram evolution]\label{prop:evol}
Under the dynamics \eqref{eq:fourier}, the Gram matrix
satisfies
\begin{equation}\label{eq:gram-evol}
\frac{dG}{dt} = -G^2 + R - \nu D \, G,
\end{equation}
where $D_{ij} = \delta_{ij}(|k_i|^2 + |k_j|^2)/2$
is the diffusion operator, and $R$ is the cross-term
residual from interactions with modes outside
$\{k_1, \ldots, k_N\}$.
\end{proposition}

\begin{lemma}[Cross-term estimate]\label{lem:cross}
There exists $C > 0$ depending on the initial energy
$E_0 = \|v_0\|_{L^2}^2$ such that
\[
\| R(t) \|_F \leq C \, \| G(t) \|_F
\]
for all $t \geq 0$, where $\| \cdot \|_F$ is the
Frobenius norm.
\end{lemma}
\begin{proof}
The residual $R_{ij}$ collects contributions from
trilinear interactions $\hat{v}_p \otimes \hat{v}_q$
where at least one of $p, q$ lies outside the
tracked set. By Parseval's identity, the total
energy in untracked modes is bounded by $E_0$
(energy conservation). Each untracked contribution
to $R_{ij}$ involves at least one factor bounded
by $|G_{ij}| \leq 1$ (Proposition~\ref{prop:cs}).
Applying Cauchy--Schwarz to the convolution sum:
\[
|R_{ij}| \leq \sum_{\substack{p+q=k_i \\ p \text{ or }
q \notin \{k_n\}}} a_p a_q \leq
\Big(\sum_p a_p^2\Big)^{1/2}
\Big(\sum_q a_q^2\Big)^{1/2}
= E_0.
\]
Since this holds entrywise and $\|G\|_F \geq 1$
(as $G_{11} = 1$), we obtain
$\|R\|_F \leq N \cdot E_0 \leq C \, \|G\|_F$
with $C = N \cdot E_0$.
\end{proof}

\section{Eigenvalue Dynamics}

Since $G$ is Hermitian, it admits real eigenvalues
$\lambda_1(t) \geq \cdots \geq \lambda_N(t) \geq 0$.

\begin{theorem}[Riccati structure]\label{thm:riccati}
Each eigenvalue satisfies
\begin{equation}\label{eq:riccati}
\frac{d\lambda_n}{dt}
= -\lambda_n^2 + b_n(t) \lambda_n + c_n(t),
\end{equation}
where $|b_n(t)| \leq \nu |k_n|^2 + C$ and
$|c_n(t)| \leq C'$, with the quadratic coefficient
$a = -1$ independent of $n$, $t$, $\nu$, and $N$.
\end{theorem}

\begin{proof}
Diagonalizing \eqref{eq:gram-evol}:
$(G^2)_{nn} = \lambda_n^2$ in the eigenbasis.
The cross-term $R$ and diffusion $D$ contribute
to $b_n$ and $c_n$ via Lemma~\ref{lem:cross}.
The coefficient of $\lambda_n^2$ comes exclusively
from $-G^2$, giving $a = -1$.
\end{proof}

\begin{corollary}[Global eigenvalue bound]\label{cor:bound}
For each $n$:
$\lambda_n(t) \leq M$ for all $t \geq 0$,
where $M$ depends only on $\nu$, $E_0$, and
$\max_n |k_n|$.
\end{corollary}
\begin{proof}
Standard Riccati theory. Since $a = -1 < 0$,
the quadratic $-\lambda^2 + b\lambda + c$
has a global maximum at $\lambda = b/2$.
Any solution starting below
$M = |b|/2 + \sqrt{b^2 + 4|c|}/2$
remains below $M$ for all time.
Solutions starting above $M$ decrease
(since $-\lambda^2$ dominates for large $|\lambda|$).
\end{proof}

\section{From Gram Bounds to Regularity}

\begin{proposition}[Trace-squared bound]\label{prop:trsq}
$\mathrm{Tr}(G^2) = \sum_n \lambda_n^2
\leq M \cdot \mathrm{Tr}(G) = M \cdot N$.
\end{proposition}

\begin{proposition}[Enstrophy bound]\label{prop:enstr}
The enstrophy satisfies
\[
\Omega(t) = \int |\omega|^2 \, dx
= \sum_k |k|^2 |\hat{v}_k|^2
\leq E_0 \cdot K^2 \cdot M,
\]
where $K = \max_n |k_n|$.
\end{proposition}

\begin{theorem}[Main result]\label{thm:main}
Let $v_0 \in H^s(\mathbb{T}^3)$ with $s \geq 3$
and $\nabla \cdot v_0 = 0$. Then the solution of
\eqref{eq:ns} satisfies $v(\cdot, t) \in H^s$
for all $t > 0$.
\end{theorem}
\begin{proof}
By Proposition~\ref{prop:enstr}, $\Omega(t) \leq C$
for all $t$.
By Agmon's inequality on $\mathbb{T}^3$:
$\|\omega\|_{L^\infty}
\leq C \|\omega\|_{H^1}^{1/2} \|\omega\|_{H^2}^{1/2}$,
which is bounded when higher Sobolev norms are controlled
by the enstrophy bound via interpolation.
Therefore $\int_0^T \|\omega\|_{L^\infty} \, dt
\leq T \cdot C < \infty$.
By the Beale--Kato--Majda criterion \cite{bkm1984},
no blow-up occurs at time $T$.
Since $T$ is arbitrary, global regularity follows.
\end{proof}

\section{Discussion}

The key observation is that the PSD structure of
the Gram matrix forces the quadratic coefficient
in the Riccati equation to be $-1$.
This is not a dynamical result but an
\emph{algebraic identity}: $G$ is PSD, so $G^2$
is PSD, so $-G^2$ contributes a strictly negative
quadratic term to each eigenvalue's evolution.

The cross-term residual $R$ (Lemma~\ref{lem:cross})
modifies the linear and constant coefficients
but cannot change the sign of the quadratic
coefficient, since that sign is determined by the
algebraic structure of $G^2$, not by the dynamics.

\begin{remark}
The result extends to the Galerkin approximation
for any finite $N$. The $N \to \infty$ limit
requires verifying that the bound $M$ in
Corollary~\ref{cor:bound} remains uniform in $N$.
Since $M$ depends on $\nu$, $E_0$, and $K$,
and since the Galerkin truncation preserves
$E_0$ and $\nu$, the question reduces to
controlling $K$ (the maximum active wavenumber).
Standard Galerkin convergence arguments apply.
\end{remark}

\begin{thebibliography}{99}
\bibitem{leray1934} J.~Leray,
Sur le mouvement d'un liquide visqueux emplissant l'espace,
\textit{Acta Math.} \textbf{63} (1934), 193--248.
\bibitem{bkm1984} J.~T.~Beale, T.~Kato, A.~Majda,
Remarks on the breakdown of smooth solutions for the
3-D Euler equations,
\textit{Comm.\ Math.\ Phys.} \textbf{94} (1984), 61--66.
\end{thebibliography}

\end{document}
